\clearpage\pagestyle{otMain}\setlength{\otBlocEspacement}{2.5mm}\fontsize{11}{15}\selectfont\corrigesChapitre{11}{Probabilités : comprendre le hasard}
\begin{blocLivre}{cor-11-1}\begin{corrige}{1}{Un dé équilibré}\demarche{Identifier les issues équiprobables.}
Les issues sont 1, 2, 3, 4, 5 et 6. Elles sont équiprobables puisque le dé est équilibré. Une seule est favorable : $\res{P(4)=1/6}$.
\end{corrige}\end{blocLivre}
\begin{blocLivre}{cor-11-2}\begin{corrige}{2}{Un événement}\demarche{Traduire précisément la condition.}
Les issues favorables sont 5 et 6, soit deux sur six. La probabilité vaut $2/6=\res{1/3}$. « Strictement » exclut la valeur 4.
\end{corrige}\end{blocLivre}
\begin{blocLivre}{cor-11-3}\begin{corrige}{3}{Certain et impossible}\demarche{Comparer l’événement à la liste des issues.}
« Obtenir un nombre inférieur à 7 » est certain : probabilité $\res1$. « Obtenir 8 » est impossible : probabilité $\res0$. Ces valeurs encadrent toutes les probabilités.
\end{corrige}\end{blocLivre}
\begin{blocLivre}{cor-11-4}\begin{corrige}{4}{Plusieurs écritures}\demarche{Passer de la fraction au quotient puis au pourcentage.}
$3/8=\res{0{,}375}$, soit $\res{37{,}5\,\%}$. Les trois écritures représentent la même probabilité.
\end{corrige}\end{blocLivre}
\begin{blocLivre}{cor-11-5}\begin{corrige}{5}{Le contraire}\demarche{Vérifier que les événements sont bien contraires.}
Perdre est ici l’événement contraire de gagner. Sa probabilité est $1-0{,}28=\res{0{,}72}$, soit 72\,\%. L’absence de résultat nul est nécessaire pour cette conclusion.
\end{corrige}\end{blocLivre}
\begin{blocLivre}{cor-11-6}\begin{corrige}{6}{Une urne}\demarche{Compter les objets équiprobables.}
Le total vaut 8. $P(R)=4/8=\res{1/2}$. Le contraire contient les 3 bleues et la verte : $P(\text{non R})=4/8=\res{1/2}$. Compter trois couleurs ne donnerait pas trois chances égales.
\label{fig-cor-11-6}\figureGeo[corrige]{urne-431}
\end{corrige}\end{blocLivre}
\begin{blocLivre}{cor-11-7}\begin{corrige}{7}{Une roue}\demarche{Repérer les secteurs favorables.}
Les numéros favorables sont 3 et 6. Deux secteurs sur huit conviennent : $2/8=\res{1/4}$. L’égalité des chances est donnée par le modèle.
\label{fig-cor-11-7}\figureGeo[corrige]{roue-8}
\end{corrige}\end{blocLivre}
\begin{blocLivre}{cor-11-8}\begin{corrige}{8}{Des faces inégales}\demarche{Utiliser la somme des probabilités.}
Le total des trois probabilités vaut 1, donc $p=1-0{,}2-0{,}5=\res{0{,}3}$. Le contraire de B a pour probabilité $1-0{,}5=\res{0{,}5}$. On ne compte pas simplement un tiers par issue.
\end{corrige}\end{blocLivre}
\begin{blocLivre}{cor-11-9}\begin{corrige}{9}{Une cible carrée}\demarche{Comparer les aires dans le modèle uniforme.}
L’aire totale est 100 cm$^2$ et l’aire gagnante 16 cm$^2$. Avec le choix uniforme, la probabilité vaut $16/100=\res{0{,}16}$, soit 16\,\%. Les longueurs seules ne sont pas le bon rapport.
\label{fig-cor-11-9}\figureGeo[corrige]{cible-carree}
\end{corrige}\end{blocLivre}
\begin{blocLivre}{cor-11-10}\begin{corrige}{10}{Une observation}\demarche{Séparer observation et modèle.}
La fréquence observée est $31/50=\res{0{,}62}$. Le modèle annoncé reste une probabilité de pile de $\res{1/2}$ au prochain lancer indépendant. Une fréquence sur un nombre fini d’essais peut s’écarter de la probabilité.
\end{corrige}\end{blocLivre}
\begin{blocLivre}{cor-11-11}\begin{corrige}{11}{Un effectif inconnu}\demarche{Retrouver le total à partir de la part rouge.}
Rouge représente un quart du total, donc le sac contient $5\times4=20$ boules. Les bleues sont $20-5=\res{15}$. Contrôle : $5/(5+15)=1/4$.
\end{corrige}\end{blocLivre}
\begin{blocLivre}{cor-11-12}\begin{corrige}{12}{Une couronne gagnante}\demarche{Soustraire les aires avant de prendre le rapport.}
L’aire totale est $36\pi$ et le disque intérieur $9\pi$. La couronne a pour aire $27\pi$. La probabilité vaut $27\pi/(36\pi)=\res{3/4}$, soit 75\,\%. On peut aussi prendre le contraire du disque intérieur.
\label{fig-cor-11-12}\figureGeo[corrige]{cible-couronne}
\end{corrige}\end{blocLivre}
\begin{blocLivre}{cor-11-13}\begin{corrige}{13}{Changer la composition}\demarche{Modifier à la fois le favorable et le total.}
Pour que les deux couleurs soient équiprobables, il faut autant de rouges que de bleues. Les bleues restent 5 ; il faut donc ajouter $\res2$ rouges. Le total augmente aussi : $5/(5+5)=1/2$.
\end{corrige}\end{blocLivre}
\begin{blocLivre}{cor-11-14}\begin{corrige}{14}{Une roue non uniforme}\demarche{Prendre la part de l’angle total.}
La probabilité du secteur de $150^\circ$ vaut $150/360=\res{5/12}$. Son contraire vaut $1-5/12=\res{7/12}$. Les trois secteurs n’ont pas la même probabilité car leurs angles diffèrent.
\label{fig-cor-11-14}\figureGeo[corrige]{roue-inegale}
\end{corrige}\end{blocLivre}
\begin{blocLivre}{cor-11-15}\begin{corrige}{15}{Des numéros et des couleurs}\demarche{Énumérer les issues favorables sans répétition.}
Les jetons favorables sont R1, R2, R3, B1 et V1 : $\res{5/6}$. Rouge a pour probabilité $3/6$ et numéro 1 aussi, mais R1 appartient aux deux événements. Leur addition le compterait deux fois. L’énumération des jetons évite ce doublon.
\label{fig-cor-11-15}\figureGeo[corrige]{jetons-numerotes}
\end{corrige}\end{blocLivre}
